\SourceLeaf{221}%
% A544-P0221-D01
eritque
\[
\begin{array}{lll}
A\equiv-(5\ast0),&B\equiv(5\ast0)+(5\ast3),&C\equiv-2-(5\ast4)\\
A'\equiv-(5\ast1),&B'\equiv(5\ast1)+(5\ast4),&C'\equiv-2-(5\ast0)\\
\text{etc.}&\text{etc.}&\text{etc.}
\end{array}
\]

% A544-P0221-D02
Quare
\[
\begin{array}{lll}
(15\ast0),\ (15\ast5),\ (15\ast10)&\begin{gathered}\text{erunt radices}\\\text{congr.}\end{gathered}&x^3-17x^2+108x-60\equiv0\\
(15\ast1),\ (15\ast6),\ (15\ast11)&&\\
(15\ast2),\ (15\ast7),\ (15\ast12)&&\\
(15\ast3),\ (15\ast8),\ (15\ast13)&&\\
(15\ast4),\ (15\ast9),\ (15\ast14)&&
\end{array}
\]

% A544-P0221-D03
Hic autem habetur
\[
\begin{gathered}
(15\ast0)^3-3(15\ast0)\equiv(15\ast1)\\
(15\ast1)^3-3(15\ast1)\equiv(15\ast2)\\
\text{etc.}
\end{gathered}
\]

% A544-P0221-D04
Unde si una radicum primae congruentiae, $10$, ponatur $(15\ast0)$ habetur
\[
\begin{array}{rlrlrl}
(15\ast0)\equiv&10&(15\ast5)\equiv&& (15\ast10)\equiv&\\
(15\ast1)\equiv&37&(15\ast6)\equiv&& (15\ast11)\equiv&\\
(15\ast2)\equiv&-151&(15\ast7)\equiv&& (15\ast12)\equiv&\\
(15\ast3)\equiv&-39&(15\ast8)\equiv&& (15\ast13)\equiv&\\
(15\ast4)\equiv&-112&(15\ast9)\equiv&& (15\ast14)\equiv&
\end{array}
\]

% A544-P0221-D05
Tandem harum singularum periodorum capiantur termini constituentes eruntque
\[
\begin{array}{lll}
(1),\ (30)&\text{radices congr.}&x^2-(15\ast0)x+1\equiv0\\
(3),\ (28)&&x^2-(15\ast1)x+1\equiv0\\
&\text{etc.}
\end{array}
\]

% A544-P0221-D06
Prima congruentiae radices sunt $126$ et $195$, quae igitur erunt radices primitivae congruentiae $x^{31}\equiv1$ et ex his reliquae sine negotio deduci possunt.
\BookRule{.28\linewidth}

\par\bigskip\Needspace{9\baselineskip}
\SourceLeaf{222}%
% A544-P0222-D01
\renewcommand{\BookRunningTitle}{Nachlass. Disquisitiones generales de congruentiis}
\pdfbookmark[0]{Disquisitiones generales de congruentiis}{disquisitiones-generales-de-congruentiis}
\begin{center}
{\large DISQUISITIONES GENERALES DE CONGRUENTIIS.}\par\medskip
ANALYSIS RESIDUORUM CAPUT OCTAVUM.
\end{center}
\BookRule{.18\linewidth}
\renewcommand{\Article}[1]{\par\medskip\Needspace{4\baselineskip}\pdfbookmark[1]{Art. #1}{disquisitiones-generales-art-#1}\centerline{\textbf{#1.}}\nobreak\smallskip}
\Article{330}
Quae in Sectionibus praecedentibus de congruentiis sunt tradita, simplicissimos tantum casus attinent methodisque particularibus plerumque sunt eruta. In hac Sectione periculum faciemus congruentiarum theoriam, quantum quidem adhuc licet, ad altiora principia reducere, simili fere modo ut \emph{aequationum} theoria considerari solet, quacum insignis intercedit analogia, uti iam saepius observavimus. Quoniam igitur omnes congruentiae algebraicae unicam incognitam involventes ad hanc formam reduci possunt
\[
X\equiv0
\]
ubi $X$ est functio algebraica incognitae $x$, nullas fractiones involvens, huiusmodi functiones imprimis erunt considerandae.

% A544-P0222-D02
\Article{331}
Si $P,Q$ sint functiones indeterminatae $x$ huius formae
\[
\begin{gathered}
A+Bx+Cxx+Dx^3+\ldots\\
H+Ix+Kxx+Lx^3+\ldots
\end{gathered}
\]
(quales abhinc semper per \emph{functiones} simpliciter designamus) et in utraque coëfficientes similium ipsius $x$ potestatum secundum quemcunque modulum sint 
\SourceLeaf{223}%
% A544-P0223-D01
congrui, \emph{functiones secundum hunc modulum congruae} dicentur. Perspicuum autem est, functiones congruas, si pro indeterminata valores aequales aut congrui accipiantur, valores congruos nancisci. Quae in Capp \textsc{i}. et \textsc{ii}. de \emph{numeris} demonstravimus, plerumque etiam de functionibus sunt tenenda; ita si $P\equiv P'$, $Q\equiv Q'$, $R\equiv R'$ etc., patet, fore $P+Q+R\text{ etc.}\equiv P'+Q'+R'+\text{etc.}$; $P-Q\equiv P'-Q'$; $PQ\equiv P'Q'$; $PQR\text{ etc.}\equiv P'Q'R'\text{ etc.}$ Demonstrationes facillimae, possuntque simili modo adornari ut Cap.\ $\mathrm{I}^{\mathrm{mo}}$.

% A544-P0223-D02
Si $PQ\equiv R$, functionem $Q$ per $\frac RP$ designabimus apposito modulo, dicemusque, $Q$ esse quotientem, si $R$ per $P$ secundum hunc modulum dividatur. Manifestum autem est, loco ipsius $Q$ omnes functiones ipsi congruas accipi posse, quas omnes tamquam \emph{unicum} valorem spectabimus. Infra vero ostendemus, quibus casibus talis quotiens plures valores (i.\,e. incongruos) nancisci possit.

% A544-P0223-D03
\Article{332}
Si modulus sit numerus primus et divisor $Q$ unicum tantum terminum involvat $Hx^h$, cuius coëfficiens $H$ per modulum non dividitur, i.\,e. si modo $H$ non sit $\equiv0$, quotiens plures valores habere nequit. Si enim esset $QA\equiv P$ et $QB\equiv P$, foret $Q(A-B)\equiv0$. Iam sit
\[
Q\equiv\dotstwo+Hx^h+Ix^{h+1}+\text{etc.}
\]
ita ut $H$ per $p$ non dividatur, et
\[
A-B\equiv Lx^l+Mx^{l+1}+\text{etc.}
\]
ita ut $L$ per $p$ non dividatur (hanc autem formam $A-B$ habebit, quia supponimus $A$ non $\equiv B$). Foretque $Q(A-B)\equiv HLx^{h+l}+\text{etc.}\equiv0$. Q.\,E.\,A., quia $HL$ non $\equiv0$.

% A544-P0223-D04
Facile iam regulae dantur functionem $P$ per $Q$, siquidem fieri potest, dividendi; sit
\[
\begin{aligned}
P&\equiv ax^\alpha+bx^{\alpha+1}+cx^{\alpha+2}+\text{etc.}+kx^\kappa\\
Q&\equiv mx^\mu+nx^{\mu+1}+qx^{\mu+2}+\text{etc.}+tx^\tau
\end{aligned}
\]
ita ut $a,k,m,t$ per modulum non dividantur, debetque esse $\alpha$ non $<\mu$, $\kappa$ non $<\tau$. Divisio autem simili modo institui potest, ut in calculo logistico communi, modo semper pro quotiente numerus integer accipiatur; scilicet quotiens semper
\SourceLeaf{224}%
% A544-P0224-D01
hanc formam habebit $\frac rm$, quod secundum modulum determinari debet. Iam si postquam $\kappa+\mu-\alpha-\tau+1$ termini sunt inventi, residuum remaneat, quod erit formae
\[
Ax^{\kappa+\mu-\tau+1}+Bx^{\kappa+\mu-\tau+2}+\dotstwo+Cx^\kappa
\]
neque omnes coëfficientes $A,B,C\dotstwo$ sint $\equiv0$, $P$ per $Q$ dividi nequit.

% A544-P0224-D02
Ceterum patet, divisionem etiam a terminis, qui maximas dimensiones habent, $kx^\kappa$, $tx^\tau$ incipi potuisse; operatio facilitabitur, si $Q$ ad formam redigatur
\[
mx^\mu(1+qx+rxx+\text{etc.})
\]
unde fiet posito $mv\equiv1$
\[
\frac PQ\equiv\frac{vP:x^\mu}{1+qx+\text{etc.}}
\]
tunc vero divisio per methodos communes perfici potest.

% A544-P0224-D03
\Article{333}
\textsc{Theorema.} \emph{Si $x\equiv a$ fuerit radix congruentiae $\xi\equiv0$, $\xi$ per $x-a$ dividi poterit secundum congruentiae modulum.}

% A544-P0224-D04
\emph{Demonstratio.} Si enim dividi non posset, foret $\xi\equiv(x-a)\xi'+b$, ita ut $b$ per modulum dividi non posset. Iam si $x$ ponatur $\equiv a$, $\xi$ fiet $\equiv0$ (hyp.), quare $(x-a)\xi'+b\equiv0$; sed tunc etiam $(x-a)\xi'\equiv0$, quare $b$ necessario erit $\equiv0$.

% A544-P0224-D05
\Article{334}
\textsc{Problema.} \emph{Datis binis functionibus, earum communem divisorem (maximae dimensionis) invenire secundum modulum datum.}

% A544-P0224-D06
\emph{Solutio.} Sint functiones $A,B$. Habeat $A$ totidem aut plures dimensiones quam $B$; dividatur $A$ per $B$, si fieri potest sine residuo, $B$ erit divisor communis quaesitus. Si residuum maneat $C$, hoc inferiorem dimensionem habebit, quam $B$. Sit itaque
\[
A\equiv aB+C,\quad B\equiv bC+D,\quad C\equiv cD+E,\quad\text{etc.}
\]
ita ut $A,B,C,D,a,b,c$ etc. sint functiones, et dimensiones functionum $A,B,C,D$ etc. constituant seriem decrescentem. Iam si tandem aliqua divisio succedat, ex.\ gr.\ $D\equiv dE$, ultimus divisor erit divisor communis quaesitus; si vero nulla succedat, tandem ad residuum pervenietur, quod nullam dimensionem
\SourceLeaf{225}%
% A544-P0225-D01
habeat i.\,e. ad \emph{numerum}; hoc autem casu functiones $A,B$ communem divisorem non habent.

% A544-P0225-D02
\emph{Demonstr.} Si divisor $E$ functionem praecedentem sine residuo dividat, omnes antecedentes dividere facile perspicitur; quare $E$ erit divisor communis functionum $A,B$. Q.\,E.\,Pr. Si autem daretur divisor maioris dimensionis, puta $E'$, hic propter $C\equiv A-aB$ etiam $C$ similique argumento etiam $D$ etc. adeoque $E$ divideret, functio maioris dimensionis functionem minoris. Q.\,E.\,A.\ Q.\,E.\,Scd. Hinc etiam patet, si divisor communis ullius dimensionis datur, ad residuum nullius dimensionis perveniri non posse; alias enim functio nullius dimensionis per functionem alicuius dimensionis divideretur. Q.\,E.\,A.

% A544-P0225-D03
\Article{335}
\textsc{Theorema.} \emph{Si $A,B$ sint functiones inter se primae secundum modulum $p$; $A$ autem dimensionis $\alpha$, $B$ dimensionis $\beta$; inveniri poterunt functiones $P,Q$, dimensionum quae sunt respective $<\beta$, $<\alpha$, ita ut}
\[
PA+QB\equiv1\;(\mathit{mod.}\,p)
\]

% A544-P0225-D04
\emph{Demonstr.} Hoc enim casu erit
\[
A\equiv aB+C,\quad B\equiv bC+D,\quad\text{etc.}\quad K\equiv kL+M
\]
ita ut dimensiones functionum $A,B,C,D,\dotstwo K,L,M$ continuo decrescant et $M$ nullam dimensionem habeat. Iam formentur series
\[
\begin{gathered}
a,\quad a',\quad a'',\quad a''',\ldots a^{(\kappa)}\\
1,\quad b,\quad b',\quad b'',\ldots b^{(\kappa-1)}
\end{gathered}
\]
ita ut
\[
\begin{array}{llll}
a'\equiv ba+1&a''\equiv ca'+a&a'''\equiv da''+a'&\text{etc.}\\
b'\equiv cb+1&b''\equiv db'+b&b'''\equiv eb''+b'&\text{etc.}
\end{array}
\]
eritque
\[
A-aB\equiv+C,\quad bA-a'B\equiv-D,\quad b'A-a''B\equiv+E,\quad\text{etc.}
\]
uti sine negotio perspicitur; hinc tandem
\[
b^{(\kappa-1)}A-a^{(\kappa)}B\equiv\pm M
\]
Iam sit $\frac1{\pm M}\equiv\mu$, eritque ponendo $P\equiv\mu b^{(\kappa-1)}$, $Q\equiv-\mu a^{(\kappa)}$
